% captions_main.tex: generated by gmpe-bench; graphics paths are relative to results/paper/

\begin{figure}[htbp]
\centering
\adjustimage{max width=\textwidth}{figures/fig1_data_splits.pdf}
\caption{Dataset and benchmark splits (full profile: 20,926 records, 575 events, 4,026 stations). (a) Magnitude versus rupture distance by hypocentre region (record counts in the legend). (b) Station $V_{S30}$; dotted lines mark the NEHRP site-class boundaries E|D at 180 m/s, D|C at 360 m/s, C|B at 760 m/s and B|A at 1500 m/s (only boundaries within the data range are drawn). (c) Usable records per period after masking $T > 1/\mathrm{LUF}$. (d) Share of test records whose event or station also occurs in training, averaged over the folds of each split scheme. Split schemes: Rand = record-level random split (leaky reference); Evt = event-grouped 5-fold; Stn = station-grouped 5-fold; Reg = leave-one-region-out; Temp = temporal (train before 2008, test 2008 onward); $M_w{\geq}7$ = magnitude hold-out (test events with $M_w \geq 7$); CL1$\rightarrow$2 = mainshocks (CL1) for training, aftershocks (CL2) for testing.}
\label{fig:fig1_data_splits}
\end{figure}

\begin{figure}[htbp]
\centering
\adjustimage{max width=\textwidth}{figures/fig2_variance_components.pdf}
\caption{Period-dependent variance components of test residuals: (a) between-event $\tau(T)$, (b) site-to-site $\phi_{S2S}(T)$ and (c) single-station within-event $\phi_0(T)$, estimated by crossed random effects on held-out events (Event split, full profile), averaged over folds.}
\label{fig:fig2_variance_components}
\end{figure}

\begin{figure}[htbp]
\centering
\adjustimage{max width=\textwidth}{figures/fig3_regime_accuracy.pdf}
\caption{Accuracy across distribution-shift regimes (full profile): (a) test RMSE and (b) empirical coverage of nominal 90\% Gaussian intervals (dashed line), averaged over folds and 27 intensity measures, for classical NGA-West2 models (train\_fit bias term), the mixed-effects baseline and machine-learning models. Hollow markers (site) are the NGA-West2 models with train\_fit station terms added (zero for stations unseen in train\_fit) and the matching single-station sigma. Split schemes: Evt = event-grouped 5-fold; Stn = station-grouped 5-fold; Reg = leave-one-region-out; Temp = temporal (train before 2008, test 2008 onward); $M_w{\geq}7$ = magnitude hold-out (test events with $M_w \geq 7$); CL1$\rightarrow$2 = mainshocks (CL1) for training, aftershocks (CL2) for testing.}
\label{fig:fig3_regime_accuracy}
\end{figure}

\begin{figure}[htbp]
\centering
\adjustimage{max width=\textwidth}{figures/fig4_conformal_ablation.pdf}
\caption{Uncertainty-method ablation at nominal 90\% (full profile): (a) empirical coverage (dashed line: nominal) and (b) mean interval width, averaged over folds and 27 intensity measures. Markers are the mean over 8 models; bars span the model range. The dashed line in (a) is the nominal 90\% target. Split schemes: Evt = event-grouped 5-fold; Stn = station-grouped 5-fold; Reg = leave-one-region-out; Temp = temporal (train before 2008, test 2008 onward); $M_w{\geq}7$ = magnitude hold-out (test events with $M_w \geq 7$); CL1$\rightarrow$2 = mainshocks (CL1) for training, aftershocks (CL2) for testing.}
\label{fig:fig4_conformal_ablation}
\end{figure}

\begin{figure}[htbp]
\centering
\adjustimage{max width=\textwidth}{figures/fig5_spectral_bands.pdf}
\caption{Example test response spectra (RotD50) with nominal 90\% bands from Hybrid ASK14+MLP (full profile; rows: Event split, fold 2; Region split, held-out region California WUS). Shaded: HSCP marginal band; dotted: split conformal band; dash-dot: HSCP simultaneous band over all usable periods; dashed: model median (HSCP re-centres records at stations with a calibration-set site term). Records are chosen at the 10th, 50th and 90th magnitude percentiles among test records with at least 90\% usable periods.}
\label{fig:fig5_spectral_bands}
\end{figure}

\begin{figure}[htbp]
\centering
\adjustimage{max width=\textwidth}{figures/fig6_conditional_coverage.pdf}
\caption{Conditional coverage at nominal 90\% (dashed) by (a) magnitude, (b) rupture distance (km), (c) NEHRP $V_{S30}$ class, ordered from stiff rock to soft soil (A/B: $\geq$760 m/s; C: 360–760 m/s; D: 180–360 m/s; E: <180 m/s), and (d) whether the test station was seen in training or calibration (full profile), pooled over split schemes, models, folds and intensity measures.}
\label{fig:fig6_conditional_coverage}
\end{figure}
